\subsection{El Enfoque Basado en Riesgo}

Los estándares internacionales en materia de prevención del lavado de dinero han desplazado su recomendación desde un modelo de cumplimiento basado en reglas fijas hacia un modelo articulado alrededor del riesgo \cite{FATF2012InternationalRecommendations, BaselCommitteeonBankingSupervision2020SoundTerrorism}. Bajo este principio, conocido como enfoque basado en riesgo, cada institución financiera debe evaluar de manera continua la exposición que representan sus clientes, sus productos, sus canales de distribución y las jurisdicciones con las que opera. A partir de esa evaluación, la institución asigna sus recursos de vigilancia de forma proporcional, de modo que las relaciones de mayor riesgo reciben controles reforzados mientras que las de menor riesgo se someten a medidas simplificadas. De esta manera, el esfuerzo de cumplimiento deja de distribuirse de forma homogénea y pasa a concentrarse donde la amenaza es más plausible.

La lógica que sustenta este enfoque descansa en la eficiencia y en la efectividad del sistema de control en su conjunto. Un régimen de reglas rígidas, que exige el mismo procedimiento a todos los clientes con independencia de su perfil, tiende a saturar a los equipos de cumplimiento con revisiones de bajo valor y a generar un volumen elevado de alertas poco informativas. En contraste, al graduar la intensidad de los controles según el riesgo, la institución libera capacidad analítica para dedicarla a los casos que verdaderamente la ameritan y mejora su probabilidad de detectar operaciones sospechosas. Además, este esquema reconoce que la actividad delictiva evoluciona, por lo que ofrece la flexibilidad necesaria para reasignar recursos cuando aparecen nuevas amenazas sin necesidad de reformar por completo el marco normativo.

Sin embargo, la aplicación práctica del enfoque basado en riesgo plantea desafíos considerables que no deben subestimarse. En primer lugar, la calidad de todo el sistema depende de la solidez de la evaluación de riesgo inicial, la cual exige información confiable, criterios coherentes y personal capacitado para interpretarla, condiciones que no siempre concurren. En segundo lugar, la discrecionalidad que otorga este modelo puede convertirse en una vía para relajar controles de forma indebida, ya sea por presiones comerciales o por una lectura interesada del riesgo, lo que obliga a las autoridades a supervisar cómo se ejerce ese juicio. Por último, la naturaleza cambiante de los productos financieros, y en particular la irrupción de nuevos instrumentos digitales, dificulta clasificar el riesgo de manera estable, de modo que la evaluación requiere una actualización permanente para no quedar desfasada.

\subsection{Desafíos de la Cooperación Transfronteriza}

El lavado de dinero constituye, por su propia naturaleza, un fenómeno que desborda las fronteras de una sola jurisdicción y se despliega a través de múltiples países. Los flujos ilícitos suelen fragmentarse y desplazarse entre distintos sistemas financieros precisamente para dificultar su rastreo y para diluir el vínculo con el delito que los origina. Cuando el circuito de ocultamiento involucra criptomonedas, esta dimensión transnacional se acentúa todavía más, dado que las redes de activos virtuales operan sobre infraestructuras globales que no reconocen los límites políticos tradicionales. En consecuencia, una investigación que se limite al territorio de un solo Estado difícilmente logra reconstruir la totalidad de la cadena de transacciones.

Frente a este carácter global, la respuesta institucional permanece anclada en gran medida a la soberanía de cada país, lo que genera fricciones importantes para la cooperación. El intercambio de información entre autoridades de distintas naciones tropieza con procedimientos formales lentos, con requisitos legales heterogéneos y con celos legítimos en torno a la protección de datos y la confidencialidad. A ello se suma que los marcos legales difieren en la manera de tipificar los delitos, en las facultades que otorgan a los organismos de supervisión y en el grado de exigencia de sus controles preventivos. Estas divergencias hacen que una operación considerada sospechosa en un país pueda pasar inadvertida en otro, y que la evidencia recolectada en una jurisdicción no siempre resulte admisible o útil en la siguiente.

Los actores dedicados al lavado de dinero conocen bien estas asimetrías y las explotan de manera deliberada mediante una práctica que suele denominarse arbitraje regulatorio. En esencia, dirigen sus operaciones hacia aquellas jurisdicciones cuyos controles resultan más débiles, cuya supervisión es más laxa o cuya cooperación internacional es más limitada, con el fin de aprovechar los eslabones más frágiles de la cadena global. De este modo, un sistema de control es tan fuerte como lo permita su componente más permisivo, y los esfuerzos rigurosos de una jurisdicción pueden verse neutralizados por la existencia de refugios normativos en otra. Por esta razón, la comunidad internacional ha insistido en la armonización de estándares y en el fortalecimiento de los mecanismos de asistencia mutua, aun cuando el avance en esa dirección continúe siendo desigual y lento.

\subsection{Proveedores de Servicios de Activos Virtuales y sus Obligaciones}

Los denominados Proveedores de Servicios de Activos Virtuales (\textit{Virtual Asset Service Providers}, VASP) constituyen las entidades que ofrecen servicios de intercambio, custodia, transferencia o administración de criptoactivos por cuenta de terceros. Dentro de esta categoría se ubican, de forma destacada, los \textit{exchanges} de criptomonedas, es decir, las plataformas que permiten convertir moneda tradicional en activos virtuales y viceversa, así como intercambiar unos criptoactivos por otros. Estos proveedores ocupan una posición estratégica en el ecosistema, pues actúan como puntos de conexión entre la economía tradicional y las redes descentralizadas. Justamente por concentrar ese tránsito, se han convertido en un eslabón crítico tanto para el funcionamiento legítimo del sector como para los intentos de canalizar fondos de origen ilícito.

Consciente de esta relevancia, la regulación reciente ha incorporado a los VASP dentro del perímetro de cumplimiento en materia de prevención del lavado de dinero, sujetándolos a obligaciones análogas a las que tradicionalmente han recaído sobre los bancos. Entre esas obligaciones figuran la identificación y verificación de la identidad de sus usuarios, conocida por sus siglas en inglés como \textit{Know Your Customer} (KYC), el monitoreo continuo de las operaciones y el reporte de aquellas que resulten sospechosas ante las autoridades competentes. A ello se añade el deber de transmitir los datos que identifican al emisor y al receptor de una transferencia de activos virtuales, un requisito que replica la lógica de la información que acompaña a las transferencias bancarias convencionales \cite{FATF2012InternationalRecommendations}. Con estas exigencias, los organismos internacionales buscan cerrar la brecha que durante un tiempo permitió que el sector cripto operara al margen de los controles aplicables al resto del sistema financiero.

No obstante, trasladar estas obligaciones a un entorno descentralizado y seudónimo genera retos que carecen de equivalente directo en la banca tradicional. La debida diligencia sobre el cliente se complica cuando las direcciones que operan en una red pública no revelan por sí mismas la identidad de las personas que están detrás de ellas, de modo que un mismo actor puede controlar numerosas direcciones sin exponer un vínculo evidente. La transmisión de los datos del emisor y del receptor resulta particularmente ardua cuando una transferencia involucra billeteras autocustodiadas o proveedores situados en jurisdicciones que no imponen requisitos equivalentes, lo que fragmenta la trazabilidad de la operación. Por añadidura, la velocidad con que surgen nuevos servicios y arreglos técnicos dentro del ecosistema tiende a superar la capacidad de los proveedores para ajustar sus controles, de manera que el cumplimiento efectivo exige un esfuerzo de adaptación constante y recursos analíticos considerables.

\subsection{La Brecha entre Regulación y Capacidad Tecnológica}

Existe una tensión persistente entre el ritmo acelerado con que innova el ecosistema de las criptomonedas y la velocidad, mucho más pausada, con que las instituciones y los reguladores logran adaptar sus herramientas de detección. El sector de los activos virtuales evoluciona de forma vertiginosa, con la aparición continua de nuevos productos, protocolos y formas de intermediación que reconfiguran la manera en que circula el valor. Los marcos normativos, en cambio, requieren procesos deliberativos extensos y ciclos de reforma que difícilmente pueden seguir ese ritmo. Como resultado, cuando una norma o un control finalmente entra en vigor, las prácticas que pretendía atender ya han mutado o han sido reemplazadas por otras.

Esta asimetría se manifiesta con especial claridad en las herramientas de detección que las instituciones emplean para vigilar sus operaciones. Los sistemas tradicionales suelen apoyarse en reglas fijas y umbrales predefinidos que disparan alertas cuando una transacción cumple ciertas condiciones establecidas de antemano. Ese diseño resulta razonable frente a comportamientos conocidos, pero pierde eficacia ante tipologías nuevas que no fueron previstas al momento de formular las reglas, pues quienes buscan ocultar fondos ajustan deliberadamente sus movimientos para permanecer por debajo de los criterios de alerta. Además, un esquema de reglas rígidas tiende a producir una elevada proporción de alertas que finalmente resultan infundadas, lo que satura a los equipos de análisis y desvía su atención de los casos verdaderamente relevantes.

Ante estas limitaciones, el campo de la prevención ha comenzado a explorar enfoques de detección más sofisticados que aspiran a superar la lógica de las reglas fijas. Buena parte del interés se ha orientado hacia métodos capaces de aprovechar la estructura relacional de las transacciones, es decir, la red de vínculos que se forma cuando el dinero se desplaza de un participante a otro y que revela patrones de conexión que un examen aislado de cada operación no alcanza a captar. La intuición que motiva esta dirección sostiene que el comportamiento ilícito rara vez se agota en una transacción individual y, en cambio, deja huellas en la forma en que se organizan los flujos entre múltiples actores. En este trabajo se retoma esa intuición como punto de partida, sin adentrarse por ahora en los detalles técnicos de tales enfoques, que se desarrollan en capítulos posteriores.
